\subsection{Solving the entire invariant boundary equation in the marked finite test}

\begin{statement}{D.52}{Proved}{An evaluated full class-two replacement and its complete homology continuation}
This statement uses the precise marked finite source
of the D.50--D.51 certificates. Write
\[
 \begin{gathered}
 G_* = A_5\times(S_3\times C_3),\qquad |G_*|=1080,\\
 s\longmapsto
  ((0,3,4,1,2),((1,0,2),0)),\\
 r\longmapsto
  ((4,1,2,0,3),((1,2,0),1)).
 \end{gathered}
 \tag{D.52.1}
\]
The permutations are in zero-based one-line notation.
Put \(K_*=\ker(\widehat B\to G_*)\), of free rank \(181\),
and let \(A_*=C_3\) be the central factor.
Let \(U_*=K_*^{\rm ab,(2)}\), \(V_*=U_*/2U_*\),
and \(W_*=(1-e_{A_*})V_*\), of dimension \(120\).
The full class-two exponent-four source is
\[
 \mathcal T_*=
 \widehat B/
 \overline{\Gamma_3(K_*)(K_*')^2K_*^4}.
 \tag{D.52.2}
\]

There is an explicitly evaluated multiplier-one
peripheral automorphism of \(\mathcal T_*\)
which fixes \(t=sr\), induces identity on \(G_*\),
and has nonidentity homology action \(T_*\), with
\[
          \operatorname{rank}T_*=181,\qquad
          \operatorname{rank}(T_*-1)=96,\qquad
          T_*^{\mathsf t}\Omega T_*=\Omega .
 \tag{D.52.3}
\]
Here \(T_*\) acts on \(V_*\) over \(\mathbb F_2\);
both ranks and the pairing identity are over \(\mathbb F_2\).
The matrix \(\Omega\) is the reduction modulo two of
the retained integral ribbon intersection form.
It replaces the particular finite D.50 homology
action excluded in D.51. The two matrices differ;
this is not an extension of that excluded action.
Both actions and both conjugator pairs project
to identity on their common invariant quotient
\(\mathcal T_*/[D,\widetilde A_*]\), where
\(D=\ker(\mathcal T_*\to G_*)\) and
\(\widetilde A_*\) is the actual order-three
complement specified below.
Their homology actions also project to identity
on the marked invariant quotient by
\((1-e_{A_*})U_*\).

The allowable moving parameters are evaluated
as actual peripheral conjugator parameters.
This pair space has dimension \(120\); each pair
induces the displayed derivation. The full remaining
boundary equation has \(13\) coordinates,
obtained from a \(3600\)-dimensional invariant
central module whose augmentation image has
rank \(3587\). In the program's specified
moving basis, a solved parameter is
\[
     \mathtt{0x2fc2bc437a2c2a199b32d59580fd1f}.
 \tag{D.52.4}
\]
The least significant bit is the first basis
coefficient. Its weight is \(60\). All linear
and quadratic coefficients, and the actual
central corrections, are included in the
certificate. Direct collection of the complete
\(6\,352\,056\)-letter corrected cusp word
gives identity in \(\mathcal T_*\).
Both torsion relations and generation are checked.

The new finite action has a compatible complete
two-adic homology continuation. Define
\[
 \begin{gathered}
 \mathfrak L_\infty=\ker(K_*\to U_*),\qquad
 \mathcal J_\infty^*=\widehat B/\mathfrak L_\infty,\qquad
 \mathcal J_2^*=\widehat B/\ker(K_*\to U_*/4U_*),\\
 \mathfrak L_c=
       \overline{\Gamma_3(K_*)(K_*')^2K_*^4},\\
 \mathcal Y_*=
       \widehat B/(\mathfrak L_\infty\cap\mathfrak L_c)
       =\mathcal J_\infty^*
             \times_{\mathcal J_2^*}\mathcal T_* .
 \end{gathered}
 \tag{D.52.5}
\]
One actual conjugator pair on this fibre product
retains the full evaluated \(\mathcal T_*\)-action.
The two projection kernels are
\[
 \begin{gathered}
 \ker(\mathcal Y_*\to\mathcal J_\infty^*)
          =\bigwedge\nolimits^2_{\mathbb F_2}V_*,
 \qquad \dim=16290,\\
 \ker(\mathcal Y_*\to\mathcal T_*)=4U_* .
 \end{gathered}
 \tag{D.52.6}
\]
The program verifies the common conjugator
coordinates and cusp identities through
homology exponents \(4,8,16\).

The full source kernel remains nonzero.
Its exact image in
\[
 \mathcal T_8^*=
 \widehat B/
 \overline{\Gamma_3(K_*)(K_*')^4K_*^8}
\]
is
\[
 \operatorname{im}(\mathfrak L_\infty\cap\mathfrak L_c
                          \to\mathcal T_8^*)
       =2\bigwedge\nolimits^2_{\mathbb Z/4}(U_*/4U_*).
 \tag{D.52.7}
\]
For the first two actual free graph generators
\(g_0,g_1\), the \(272\)-letter word
\([g_0,g_1]^2\) is in this source kernel and
has nonzero coordinate \(2\) at the corresponding
commutator in \(\mathcal T_8^*\).

The actual arithmetic \(G_3\), the complete
D.50--D.51 arithmetic actions, and their
remaining invariant residual are unchanged
by this calculation. No map transferring
the finite solution (D.52.4) to that arithmetic
residual is asserted. The evaluated finite
repair supplies exact allowable-correction
and boundary-equation data for that next
comparison.
The manuscript passages are Part D(e),
section 51.16, folio 734 verso, scan 174,
and section 51.15, folios 725--727,
scans 156--160.
\end{statement}
\begin{proof}
\emph{The actual parameter space and its peripheral maps.}
All group operations, the \(1080\) marked states,
the \(181\) free graph cycles, and the central
three-deck action are rebuilt from (D.52.1).
The program checks that \(1-t\) has rank
\(120\) on \(W_*\). For \(w\in W_*\), set
\[
 \begin{gathered}
 b(w)=(1-t)^{-1}s(1-r)w,\qquad a(w)=w+b(w),\\
 (1-s)b(w)+s(1-r)a(w)=0 .
 \end{gathered}
 \tag{D.52.8}
\]
Conversely every pair in the kernel of this
cusp map is of this form: take \(w=a+b\)
and solve the displayed invertible equation.
This proves that these are all \(120\)
allowable moving parameters at the first level.

The corrections give the actual derivation
values
\[
            \delta(s)=(1-s)b,\qquad
            \delta(r)=(1-r)a,\qquad
            \delta(t)=0 .
 \tag{D.52.9}
\]
They satisfy the two torsion relations and
the cusp relation as crossed homomorphisms.
Their restrictions to the free kernel are
therefore actual \(G_*\)-linear homology maps,
with zero transgression. This construction
does not replace them by all abstract
\(\operatorname{Hom}_{G_*}(V_*,W_*)\).
The edge map sends the identity edge to
itself plus \(w=a+b\), giving every matrix
column used in the calculation.

Each chosen pair lifts integrally modulo
four in \((1-e_{A_*})(U_*/4U_*)\).
The odd averaging idempotent first gives
moving lifts; the onto modulo-two cusp
map then corrects their exact modulo-four
cusp error. An actual order-three lift
\(\widetilde z\) of the central deck generator
is obtained, and checked, as in D.51.
On the moving abelian summand,
\((1-z)^{-1}=2-z\) modulo four.
Thus the two abelian coordinates lift to
literal commutators with \(\widetilde z\).
The resulting words \(A_i,B_i\) lie in
\(E=[D,\widetilde z]\), have the prescribed
homology coordinates, and have central cusp
errors. All these assertions are checked
in the faithful collected graph model.

\emph{The exact invariant quotient and quadratic boundary map.}
The central residual is
\(Z=W_*\wedge V_*\). Its invariant part
\(Z_0\) has dimension \(3600\), and its
moving part has dimension \(10860\), by D.51.
An adapted basis consists of \(61\) fixed
homology vectors and \(60\) two-dimensional
blocks \(u_i,v_i\), with
\[
                    zu_i=v_i,\qquad zv_i=u_i+v_i.
\]
A basis of \(Z_0\) is
\[
 \begin{gathered}
 u_i\wedge v_i,\\
 u_i\wedge u_j+u_i\wedge v_j+v_i\wedge v_j,\qquad
 u_i\wedge v_j+v_i\wedge u_j
                  \quad(i<j).
 \end{gathered}
 \tag{D.52.10}
\]
This lists \(60+2\binom{60}{2}=3600\)
independent invariant coordinates.
The actual exterior actions of \(s,r\)
give the full augmentation matrix
\[
       (1-s)Z_0+(1-r)Z_0\ \longrightarrow Z_0 .
 \tag{D.52.11}
\]
Its rank is \(3587\). Exact elimination
constructs \(13\) dual coordinate functionals
for its cokernel; they are retained in the
certificate and annihilate every column.
The moving central cusp equation is already
onto. Thus the remaining error is evaluated
in these actual \(13\) coordinates.

Here is the quadratic formula, including its
source in the word relation. For a pair of
kernel elements \(A,B\), write its cusp error as
\[
       C(A,B)=B\,s(B)^{-1}s(A)\,t(A)^{-1},
 \tag{D.52.12}
\]
where \(s,t\) denote their conjugation actions.
For the chosen pairs this is central, since
the abelian cusp equation was solved modulo
four. Let \(a,b\) and \(a',b'\) be their
modulo-two abelian coordinates. Collecting
the product of two pairs in the class-two
group gives
\[
 \begin{split}
 C(AA',BB')={}&C(A,B)+C(A',B')\\
             &+(1+s)b'\wedge s(a+b).
 \end{split}
 \tag{D.52.13}
\]
Central coordinates are written additively.
To verify the cross term, order the four
factors in (D.52.12) for each pair and move
the primed factors into their product-pair
positions. The cross commutators sum to
\[
   (b'+sb')\wedge(sb+sa+ta)
       +(sa'+ta')\wedge ta
       =(1+s)b'\wedge s(a+b),
\]
using the two abelian cusp equations.
This proves the formula with no linearization
of the group relation.

Choose the program's ordered moving basis
and the corresponding actual pairs \(A_i,B_i\).
For a binary parameter vector \(v\), use the
ordered products of the selected pairs.
The exact remaining obstruction is
\[
       \omega(v)=\sum_i\lambda_i v_i
                   +\sum_{i<j}\gamma_{ij}v_i v_j
                   \ \in\mathbb F_2^{13}.
 \tag{D.52.14}
\]
The coefficients are obtained by applying
the invariant projection and the \(13\)
cokernel functionals to the actual individual
errors and the cross wedges in (D.52.13).
The program stores every coefficient and
checks the polar identity again on \(36\)
literal two-parameter word products.
Thus the evaluated equations retain the
torsion, cusp and transgression constraints.

\emph{The unsuccessful pairing-only route is also evaluated.}
The earlier attempted correction solves only
the intersection-form equations. Their
linearized rank is \(9\), leaving \(111\)
directions. A selected nonidentity direction
has homology derivation rank \(16\), gives
an invertible pairing-preserving matrix, and
satisfies the exact abelian cusp equation.
Its full \(13\)-coordinate obstruction is
nevertheless nonzero.

The program also evaluates all moving
modulo-four homology variations with that
same modulo-two conjugator pair. They are
the \(120\) pairs in (D.52.8), realized by
moving projections of actual squares.
They lie in the central Frattini subgroup
of the class-two group, so their effects
are exactly linear, and are checked there.
Their additional image in the \(13\)-coordinate
cokernel has rank zero. This records the
precise failure of that attempted repair;
it is not silently discarded.

\emph{The solved parameter and full word verification.}
Evaluation of all equations (D.52.14) gives
zero at the parameter (D.52.4).
The fixed deterministic search first reaches
it at trial \(5784\). The selected product
has \(60\) basis factors. Its full cusp
word is collected directly, and its
invariant error lies in the actual
augmentation image (D.52.11).
Exact elimination supplies invariant
central corrections. The moving correction
is obtained by the literal odd-word
telescoping used in D.51.

Each central correction is realized by
its product of actual free-generator
commutator words. The program then collects
all \(6\,352\,056\) letters in the corrected
cusp word in the full class-two quotient,
and obtains zero, including the invariant
coordinates. It also checks the squares
and cubes of the two torsion conjugates.
The full homology matrix has rank \(181\)
and deviation rank \(96\); it preserves
the exact intersection form. An onto
Frattini action generates the finite
two-group kernel, so the resulting
endomorphism is an automorphism of
\(\mathcal T_*\).

All conjugators lie in \(E\), including
both central corrections. Hence their
images in \(\mathcal T_*/E\) are identity.
This is an actual quotient: the subgroup
\(E=[D,\widetilde A_*]\) is normal in \(D\);
two odd complements above \(A_*\) are
\(D\)-conjugate and give the same \(E\).
Since \(A_*\) is normal in \(G_*\), conjugation
by the full group preserves \(E\).
This proves the stated comparison to the
earlier finite actions. The new modulo-two
matrix differs from the excluded one:
the latter has pairing-defect rank \(72\),
whereas the new defect is zero. The
replacement is therefore explicit at
the common homology quotient.

\emph{The complete homology continuation and exact common source.}
The maximal common quotient of
\(\mathcal J_\infty^*\) and \(\mathcal T_*\)
is exactly \(\mathcal J_2^*\).
Indeed the image of \(\mathfrak L_c\)
in \(U_*\) is \(4U_*\): its commutator
and third lower-central generators have
zero abelian image, and its fourth powers
give all of \(4U_*\).
The image of \(\mathfrak L_\infty\) in
\(\mathcal T_*\) is exactly \(D'\):
it contains the actual commutator generators,
and its abelian image is zero.
These two original kernel computations
give (D.52.5)--(D.52.6).

The evaluated conjugator coordinates modulo
four lie in the integral moving summand.
Lift them successively through homology
exponents \(2^k\). At each step the cusp
error is in the next copy of \(W_*\).
The onto map in (D.52.8) supplies corrections.
Every integral correction is first projected
by the full odd averaging idempotent, so
it remains in the moving summand at all
later levels. For \(k\ge3\) the changes
are multiples of four and preserve the
common \(\mathcal J_2^*\) coordinates.
The unchanged invertible reduction modulo
two proves generation at every level.

The resulting automorphisms, their inverses,
and both conjugators are compatible.
Their limits define an automorphism on
\(\mathcal J_\infty^*\), whose action and
conjugators agree with the full evaluated
\(\mathcal T_*\) component on
\(\mathcal J_2^*\). The fibre product gives
one pair and one automorphism on
\(\mathcal Y_*\). The program checks
the successive \(4,8,16\) coordinates,
the exact cusp identities, and their
unchanged common coordinates.

\emph{The remaining full-source kernel has an exact finite image.}
In \(\mathcal T_8^*\), the kernel over
\(G_*\) has generator exponents modulo
eight and independent commutator coordinates
modulo four. Its class-two collection law
realizes the commutator module
\(\bigwedge^2_{\mathbb Z/4}(U_*/4U_*)\).
The image of \(\mathfrak L_c\) is generated
by fourth powers of generators and squares
of commutators. It has abelian coordinates
divisible by four and even commutator
coordinates: collecting a fourth power
adds \(\binom42=6\), which is even modulo four.

An element of
\(\mathfrak L_c\cap\mathfrak L_\infty\)
has zero abelian coordinate, so its image
is in twice that commutator module.
Conversely every commutator square has an
actual preimage in both original kernels.
This proves equality in (D.52.7).
The program evaluates \([g_0,g_1]^2\)
in both faithful normal forms. It is
identity in \(\mathcal T_*\) and has
coordinate \(2\) at \((0,1)\) in
\(\mathcal T_8^*\). Its actual \(272\)-letter
\(B\)-word is retained in the certificate.

This completes the finite replacement,
its complete homology continuation, and
its remaining kernel computation. The
arithmetic \(G_3\) has a different specified
source map, whose actual correction and
transgression comparison is still being
determined; the present proof makes no
transfer between those two marked sources.
\end{proof}
